% 本文件由 LaTeX 排版 Agent 根据有效 translation.json 自动生成。
% author=Athanasius; work_id=2820; title=无首历史

\chapter{无首历史}

% structure: p0001 source=h1 target=omitted reason=page-title-already-rendered-by-parent

I.~1. 皇帝 \Person{君士坦提乌斯} 也写信谈及 \Person{亚大纳削} 的回归，在皇帝的信函中亦可见此信。\par

2. 在 \Person{额我略} 死后，\Person{亚大纳削} 从 \Place{罗马} 城和 \Place{意大利} 地区返回，于保斐月二十四日进入 \Place{亚历山大里亚}，时值执政官 \Person{君士坦提乌斯} 四任与 \Person{康斯坦斯} 三任之年（346年10月21日）；即 [vii] 年 vi [个月又 iii 日] 之后，并在 \Place{亚历山大里亚} 安静居住了 ix 年 iii 个月 [又 xix 日]。\par

II. 他返回之后，时值执政官 \Person{利墨尼乌斯} 与 \Person{卡图利努斯} 之年（349年），\Person{德奥多勒}、\Person{纳尔基索斯} 和 \Person{乔治}，与其他一些人，来到 \Place{君士坦丁堡}，想要说服 \Person{保禄} 与他们共融，但 \Person{保禄} 甚至不愿对他们说一句话，并以其问候回应以绝罚。于是他们拉拢 \Place{尼科墨底亚} 的 \Person{欧瑟比}，为真福 \Person{保禄} 设下陷阱，并捏造关于他涉及 \Person{康斯坦斯} 和 \Person{马格嫩提乌斯} 的诽谤，将他驱逐出 \Place{君士坦丁堡}，以便他们能在那里立足，并散播 \OriginalTerm{亚略异端}{Arian heresy}。\Place{君士坦丁堡} 的民众爱戴至福的 \Person{保禄}，不断引发骚乱，以防他被带离城市，因为他们热爱他那纯正的教义。然而皇帝大怒，派伯爵 \Person{赫尔摩格内斯} 将他赶出去；但民众听到此事，将 \Person{赫尔摩格内斯} 拉到城中心论处。他们借此事件获得攻击该主教的借口，并将他流放到 \Place{亚美尼亚}。\Person{德奥多勒} 和其余的人想要在该城的 \OriginalTerm{宗座}{See} 上立一位亚略异端的盟友与同党——原 \Place{日耳曼尼西亚} 的主教 \Person{尤多克修斯}，正欲如此行时，民众被激怒引发暴动，不许任何人坐上真福 \Person{保禄} 的宗座——他们遂将 \Person{保禄} 的一位司铎 \Person{马塞多尼} 拉来，并祝圣他为 \Place{君士坦丁堡} 城的主教；全体主教集会谴责他，因为他背叛了自己的父亲，不忠地从异端者手中接受了 \OriginalTerm{覆手}{laying on of hands}。\par

然而，\Person{马塞多尼} 在与他们共融并签署后，他们便寻来无关紧要的借口，将他从教会中移除，并立了前述的 \Place{安提阿} 的 \Person{尤多克修斯}；自此，参与这次分裂的人被称为马塞多尼派，他们在 \OriginalTerm{圣神}{Holy Spirit} 方面遭遇了破船之祸。\par

III.~3. 此后，\Person{亚大纳削} 听说他将遭到骚扰，其时皇帝 \Person{君士坦提乌斯} 正住在 \Place{米兰}（353年），于是派了一艘船载着五位主教前往宫廷：\Place{特穆伊斯} 的 \Person{塞拉皮翁}、\Place{尼科塔斯} 的 \Person{特里亚德尔弗斯}、\Place{上基诺波利斯} 的 \Person{阿波罗}、\Place{帕克蒙} 的 \Person{阿摩尼乌斯}，…以及三位亚历山大里亚的司铎：医生 \Person{伯多禄}、\Person{阿斯特里库斯} 和 \Person{菲莱亚斯}。他们从 \Place{亚历山大里亚} 启航后，时值执政官君士坦提乌斯奥古斯都六任与君士坦提乌斯凯撒二任之年，于帕科蒙月二十四日（353年5月19日）；就在四天后的帕科蒙月二十八日，宫廷来的 \Person{蒙塔努斯} 进入 \Place{亚历山大里亚}，并将同一君士坦提乌斯奥古斯都的一封信交给主教 \Person{亚大纳削}，禁止他前往宫廷，因此主教极为忧伤，全体民众也深感不安。于是 \Person{蒙塔努斯} 一事无成，便离开了，留下主教在 \Place{亚历山大里亚}。\par

4. 此后不久，帝国公证官 \Person{狄奥格内斯} 于默索月（355年8月）来到 \Place{亚历山大里亚}，时值执政官 \Person{阿尔贝提翁} 与 \Person{洛利安努斯} 之年：即 \Person{蒙塔努斯} 离开亚历山大里亚后两年零五个月。\Person{狄奥格内斯} 迫切地催促每一个人，以迫使主教离开城市，并使所有人遭受不小痛苦。在托特月第六日，他进行了强烈的尝试以围困教堂，并为此努力了四个月，即从默索月，或自那闰日的[第一]天起，直到科雅克月第二十六日（12月23日）。但由于民众及法官们强烈抵抗 \Person{狄奥格内斯}，\Person{狄奥格内斯} 无功而返，即在上述科雅克月的第二十六日，执政官 \Person{阿尔贝提翁} 与 \Person{洛利安努斯} 之年，经过前述的四个月。\par

IV.~5. 在执政官 \Person{阿尔贝提翁} 与 \Person{洛利安努斯} 之后，公爵 \Person{叙利亚努斯} 和公证官 \Person{希拉里} 于提比月第十日（356年1月6日）从 \Place{埃及} 来到 \Place{亚历山大里亚}。并率领 \Place{埃及} 和 \Place{利比亚} 的所有军团兵士先行，公爵和公证官在梅基尔月第十三日，即第十四日的前夜，率全军夜间进入了 \OriginalTerm{提奥纳教堂}{Church of Theonas}。他们破开提奥纳教堂的门，带着不计其数的士兵进入其中。但主教 \Person{亚大纳削} 逃脱了他们的手，于前述梅基尔月第十四日获救。这事发生在主教自 \Place{意大利} 返回后九年三个月零十九天。但当主教获救后，他的司铎们和民众保持了对各教堂的掌控，并维持了四个月的共融，直至帕希尼月第十六日，时值执政官君士坦提乌斯八任与尤利安凯撒一任（356年6月10日），总督 \Person{卡塔弗罗尼乌斯} 和伯爵 \Person{赫拉克利乌斯} 进入 \Place{亚历山大里亚}。\par

V.~6. 他们进入后四天，亚大纳削派从各教堂中被驱逐，教堂被交给属于 \Person{乔治} 的人，他们期待 \Person{乔治} 到来为主教。因此，他们在帕希尼月第二十一日接收了各教堂。此外，\Person{乔治} 于执政官君士坦提乌斯九任与尤利安凯撒二任之年，即梅基尔月三十日（357年2月24日）抵达 \Place{亚历山大里亚}，那是在他的党羽接收教堂后八个月零十一天。于是 \Person{乔治} 进入 \Place{亚历山大里亚}，占据了各教堂整整十八个月：随后平民在狄奥尼修斯教堂袭击了他，他在托特月第一日，时值执政官 \Person{塔提安} 与 \Person{克雷阿利斯} 之年（358年8月29日），经历危险与巨大挣扎才勉强脱身。\Person{乔治} 在暴动后的第十日，即保斐月五日（10月2日）被逐出 \Place{亚历山大里亚}。但属于 \Person{亚大纳削} 主教的人，在 \Person{乔治} 离开后九天，即保[斐]月十四日，驱逐了乔治的人，并占据了各教堂两个月零十四天；直到公爵 \Person{塞巴斯蒂安} 从 \Place{埃及} 来到并驱逐了他们，并在科雅克月第二十八日（12月24日）再次将各教堂分配给乔治一派。\par

7. 自\Person{乔治}离开\Place{亚历山大城}整整九个月之后，公证人\Person{保禄}于帕伊尼月29日，\OriginalTerm{执政官}{Coss.}\ \Person{尤西比乌斯}与\Person{希帕提乌斯}任内\EditorialNote{359年6月23日}抵达，并代表乔治颁布了一道御旨，且为报复他而胁迫了许多人。而\EditorialNote{两年又}五个月后，乔治在阿提尔月30日来到亚历山大城\EditorialNote{执政官陶鲁斯与弗洛伦提乌斯}，从朝廷而来\EditorialNote{361年11月26日}，即他逃离后三年又两个月。在\Place{安提阿}，阿里乌斯异端者将保禄派逐出教会，并任命了\Person{梅勒蒂乌斯}。当他不同意他们的邪恶思想时，他们便祝圣了亚历山大城的乔治的司铎\Person{欧佐伊乌斯}来代替他。\par

VI. 8. 如前所述，\Person{乔治}于阿提尔月30日进入\Place{亚历山大城}，在城中平安住了三天，即\EditorialNote{直到}乔亚克月3日。因为在该月第4日，总督\Person{格龙提乌斯}宣布了皇帝\Person{君士坦提乌斯}驾崩的消息，且\Person{尤利安努斯}一人掌控了整个帝国。听到这消息，亚历山大城的市民和所有人一起高呼反对乔治，并一致将其拘禁。他被铁链捆绑在狱中，从上述乔亚克月4日起至同月27日止，共24天。因为同月28日清晨，几乎全城的人将乔治从狱中提出，连同与他在一起的伯爵，即被称为凯撒里昂教堂的建筑总监，将他们二人杀死，并拖着他们的尸体绕城示威，乔治的尸体被放在骆驼上，而\Person{德拉康提乌斯}的尸体则被人用绳索拖拽；如此凌辱之后，大约在当天第七时辰，他们将双方的尸体焚烧。\par

VII. 9. 在接下来的……麦基尔月10日，\OriginalTerm{执政官}{Coss.}\ \Person{陶鲁斯}与\Person{弗洛伦提乌斯}之后\EditorialNote{362年2月4日}，颁布了皇帝\Person{尤利安努斯}的一道敕令，命令将那些从前从偶像、殿宇侍者及公共账目中取走之物归还回去。\par

10. 然而三天后，即麦基尔月14日，皇帝\Person{尤利安努斯}和副主教\Person{莫德斯图斯}下达了一道命令给总督\Person{格龙提乌斯}，命令所有迄今因党争被挫败并遭流放的主教返回自己的城镇和省份。这封信于次日麦基尔月15日公布，随后总督格龙提乌斯也发布了一道谕令，命令\Person{亚大纳削}主教返回他的教会。在这道谕令颁布12天后，亚大纳削出现在\Place{亚历山大城}，并于同月麦基尔月27日进入教会；因此，从他于\Person{叙里亚努斯}和\Person{希拉里}时期逃亡起至回归时止，当尤利安努斯……麦基尔月27日。他留在教会直到帕奥菲月26日，\OriginalTerm{执政官}{Coss.}\ \Person{马梅尔提努斯}与\Person{内维塔}任内\EditorialNote{362年10月23日}，整整八个月。\par

11. 而在前述帕奥菲月27日，他\EditorialNote{总督}颁布了皇帝\Person{尤利安努斯}的一道敕令，要求\Person{亚大纳削}主教离开\Place{亚历山大城}，敕令一经颁布，主教便离开该城，住在特雷乌附近。他离开后不久，总督\Person{奥林普斯}遵从同一位\Person{皮提奥多鲁斯}以及与他在一起的极为难缠之人的命令，将亚历山大城的司铎\Person{保禄}与\Person{阿斯泰里奇乌斯}放逐，并指示他们居住在\Place{安德罗波利斯}城。\par

VIII. 12. 那位\Person{奥林普斯}总督，在门索尔月26日，\OriginalTerm{执政官}{Coss.}\ \Person{尤利安努斯·奥古斯都}四度及\Person{萨卢斯提乌斯}任内\EditorialNote{363年8月20日}，宣布皇帝\Person{尤利安努斯}已死，而\Person{约维安努斯}，一位基督徒，已为皇帝。在随后的托特月18日，皇帝约维安努斯的一封信送到奥林普斯总督处，信中要求只应崇拜至高者天主和基督，并且民众应在各教会中保持共融，践行宗教。此外，前述司铎\Person{保禄}和\Person{阿斯泰里奇乌斯}从\Place{安德罗波利斯}城的流放地返回，并在托特月10日进入\Place{亚历山大城}，历经十个月。\par

13. \Person{亚大纳削}主教如前所述在特雷翁停留后，便上到埃及的高处，直至\Place{底比斯}的\Place{赫尔莫波利斯}，甚至到了\Place{安提诺波利斯}。当他逗留在这些地方时，得知皇帝\Person{尤利安努斯}已死，而\Person{约维安努斯}，一位基督徒，已为皇帝。于是主教秘密进入\Place{亚历山大城}，他的到来不为许多人所知，并乘船沿海路去迎接皇帝约维安努斯，随后，在教会事务处理妥当后，收到了一封信，便返回亚历山大城，于\OriginalTerm{执政官}{Coss.}\ \Person{约维安努斯}与\Person{瓦罗尼安努斯}任内的阿提尔月19日进入教会。从他遵照\Person{尤利安努斯}的敕令离开亚历山大城，直至他在前述阿提尔月19日到达，经过了一年三个月又二十二天。\par

IX. 在\Place{君士坦丁堡}，\Place{日尔曼尼西亚}的\Person{欧多克修斯}掌控教会，他与\Person{马策多尼乌斯}之间存在分裂；但通过欧多克修斯，从阿里乌斯派的虚假\EditorialNote{教导}中又产生了更糟糕的异端，即\Place{尼西亚}的\Person{阿埃提乌斯}和\Person{帕特里奇乌斯}，他们与\Person{优诺米乌斯}、\Person{赫利奥多鲁斯}和\Person{斯德望}相通。欧多克修斯采纳了这一异端，并与\Place{安提阿}的阿里乌斯派主教\Person{欧佐伊乌斯}相通，他们以借口罢免了\Person{塞琉西乌斯}、\Person{马策多尼乌斯}、\Person{希帕提安}以及其他十五位属于他们的主教，因为他们不接受\Quote{不同}也不接受\Quote{受造于非受造者}。他们的阐述如下：\par

\Person{帕特里奇乌斯}和\Person{阿埃提乌斯}的阐述，他们与\Person{优诺米乌斯}、\Person{赫利奥多鲁斯}和\Person{斯德望}相通。\par

这些是\OriginalTerm{天主}{God}的属性：\OriginalTerm{不受生}{Unbegotten}、无始源、永恒、不受支配、不可变更、全视、无限、无可比拟、全能、无需预知而预知未来；无开始。这些属性并不属于\OriginalTerm{圣子}{Son}，因为祂受支配，在被支配之下，从无中被造，有终结，不与［\OriginalTerm{圣父}{Father}］相比，圣父超越祂……论到\OriginalTerm{基督}{Christ}：就涉及圣父而言，祂对将来无知。祂过去不是\OriginalTerm{天主}{God}，而是\OriginalTerm{圣子}{Son of God}；是对在祂以后的人的天主：在这一点上，祂与圣父有不变的相似性，即祂见一切事物，因为一切事物……因为祂在美善方面没有改变；［但］不是在\OriginalTerm{天主性}{Godhead}的品质上相似，也不是在\OriginalTerm{本性}{Nature}上相似。但如果我们说，祂是由天主性的品质所生，我们就是说祂类似蛇的后裔，这是不虔敬的话：如同雕像从自身生锈，并将被这锈所腐蚀，圣子也是如此，如果祂出自圣父的本性，就会腐蚀圣父。但从工程和工程的新颖来说，圣子本性上是天主，不是出自圣父的本性，而是出自另一本性，相似于圣父，但不是从祂而来。因为祂被造为天主的肖像，而我们出于天主，并源于天主。既然万物源于天主，圣子也是如此，仿佛源于某其他东西。如同铁若有锈就会减少，如同身体若生虫就会被吞噬，如同伤口若产生脓液就会被其腐蚀，同样，说圣子出自圣父本性的人［所想的］也是如此；现在，凡不说圣子相似圣父的人，当被逐出\OriginalTerm{教会}{Church}，并处以\OriginalTerm{绝罚}{anathema}。如果我们说\OriginalTerm{圣子}{Son of God}是天主，我们就引入了两个无始者：我们称祂为天主的肖像；称呼祂\Quote{出自天主}的人是\OriginalTerm{撒伯流主义化}{Sabellianises}。凡说不知天主的降生，这人就是\OriginalTerm{摩尼教色彩}{Manicheanizes}：若有人说\OriginalTerm{圣子}{Son}的\OriginalTerm{性体}{Essence}相似那不受生的圣父的性体，就是亵渎。因为如同雪和白铅在白色上相似，但在种类上不同，同样圣子的性体也与圣父的性体有别。但雪有另一种白色……\par

请乐意听闻，\OriginalTerm{圣子}{Son}在祂的作为上与\OriginalTerm{圣父}{Father}相似；正如天使无法领悟总领天神的本性，让他们乐意理解，总领天神也无法领悟\OriginalTerm{革鲁宾}{Cherubin}的本性，革鲁宾也无法领悟\OriginalTerm{圣神}{Holy Spirit}的本性，圣神也无法领悟\OriginalTerm{独生子}{Only-begotten}的本性，独生子也无法领悟\OriginalTerm{不受生的天主}{Unbegotten God}的本性。\par

14. 那时，\Person{亚大纳修}主教即将从\Place{安提阿}来到\Place{亚历山大里亚}，亚略派的\Person{欧多克修斯}、\Person{狄奥多尔}、\Person{索弗罗尼乌}、\Person{欧佐伊乌斯}和\Person{希拉里}商议，任命了\Person{乔治}的一位长老\Person{路爵乌斯}，前往\Place{皇宫}觐见\Person{约维安}皇帝，并陈明文件副本中所包含的内容。\EditorialNote{此处我们省略了一些不太必要的事项。}\par

10. 15. \Person{约维安}之后，\Person{瓦伦提尼安}和\Person{瓦伦斯}被迅速召登宝位，他们的一道法令传遍各地，也在\Place{亚历山大里亚}于帕雄月10日（即瓦伦提尼安与瓦伦斯执政官年，公元365年5月5日）颁布，大意为：那些在\Person{君士坦提乌斯}治下被废黜并逐出自己教会的主教们，他们曾在\Person{犹利安}在位期间收回并夺回了主教职位，如今应再次从教会中被逐出，各法庭将被罚三百磅黄金，除非他们将主教们逐出教会与城镇。为此，亚历山大里亚发生了大骚乱与暴动，以致整个教会陷入不安，因为当时长官\Person{弗拉维安}及其随从的人员甚少；且因皇帝敕令与罚金，他们催促主教们离开城镇；基督徒群众则抗拒反驳官员和法官，坚称\Person{亚大纳修}主教不在此规定和皇帝敕令的范围内，因为君士坦提乌斯并没有放逐他，反而是恢复了他的职位。同样，犹利安迫害他；犹利安召回了所有人，却为了偶像崇拜再次将他逐出，而约维安将他带了回来。这场对抗和骚乱持续到次月帕尼月第14日；因为在这一天，长官弗拉维安呈报，声明已就亚历山大里亚所引发的此事咨商皇帝们，于是不久众人都平息了。\par

11. 16. 四个月又二十四天之后，即帕奥菲月第8日，\Person{亚大纳修}主教夜间悄悄离开教堂，退居到靠近\Place{新河}的一座别墅。但长官\Person{弗拉维安}与军事长官\Person{维克托里努斯}不知他已退去，当夜带着一队士兵来到\Place{狄奥尼修斯堂}：他们打破后门，进入房舍上层搜寻主教的住处，却没有找到他，因为他不久前才退去。他从上述日期帕奥菲月第8日起就一直待在上述庄园，直到梅希尔月第6日，即整整四个月（10月5日至1月31日）。此后，帝国书记官\Person{布雷西达斯}在同月梅希尔月带着一封皇帝信函来到\Place{亚历山大里亚}，命令上述亚大纳修主教返回城中，照常掌管各教堂；在梅希尔月第7日，即瓦伦提尼安与瓦伦斯执政官年之后（亦即\Person{格拉提安}与\Person{德加拉伊夫斯}执政官年），上述书记官布雷西达斯随同军事长官维克托里努斯与长官弗拉维安聚集于\Place{宫殿}，向在场的法庭官员和民众宣布，皇帝们已命令主教返回城中，于是上述书记官布雷西达斯立即随同法庭官员和众多基督徒民众前往前述别墅，带着皇帝敕令迎接亚大纳修主教，在梅希尔月第7日领他进入称为\Place{狄奥尼修斯堂}的教堂。\par

十二. 17. 从执政官\Person{格拉提安}和\Person{达加拉伊夫斯}(366年)到接下来的\Person{卢皮基努斯}和\Person{约维努斯}执政官任期(367年)，以及[\Person{瓦伦提尼安二世}和]\Person{瓦伦斯二世}执政官任期，在帕伊尼月14日(368年6月8日)，在[这个]执政官任期中，\Person{亚大纳削}的[主教职]满40年。其中，他在\Place{高卢}的\Place{特里尔}住了[2年4个月11天，在\Place{意大利}和\Place{西方}]90个月零3天。在\Place{亚历山大}[和]某些隐蔽地点，当他被书记官\Person{伊拉略}和公爵骚扰时，他躲藏了72个月零14天。在\Place{埃及}和\Place{安提阿}的旅程中有15个月零22天：在新河附近的产业上住了4个月。总计正好是6个月零17年零20天。此外，他在\Place{亚历山大}安宁地住了22年零5个月10天。而且，他在最后一次旅程中及\Place{提洛}和\Place{君士坦丁堡}之外短暂停留了两次。因此，结果将如我上文所述，\Person{亚大纳削}的主教任期总共40年，直到帕伊尼月[1]4日，\Person{瓦伦提尼安}和\Person{瓦伦斯}执政官任期。在接下来的\Person{瓦伦提尼安}和\Person{维克托}执政官任期，帕伊尼月14日，1年，以及在接下来的\Person{瓦伦提尼安[三世]}和\Person{瓦伦斯三世}执政官任期，帕伊尼月14日，在接下来的\Person{格拉提安}和\Person{普罗布斯}执政官任期，[以及之后的\Person{莫德斯图斯}和\Person{阿林特乌斯}执政官任期]，以及\Person{瓦伦提尼安[四世]}和\Person{瓦伦斯四世}的另一个执政官任期，在帕孔月8日他安息(373年5月3日)。\par

十三. 18. 在上述的\Person{卢皮基努斯}和\Person{约维努斯}执政官任期内，\Person{路西乌斯}特别渴望为自己获取亚略异端的主教职位，在他离开\Place{亚历山大}很久之后，于上述执政官任期内抵达，并在托特月26日(367年9月24日)夜间秘密进城：据称，他藏匿在某座小屋里度过那一天。但第二天他去了他母亲所住的房子；他的到来立刻传遍全城，全体民众聚集起来，谴责他的进城。公爵\Person{图拉真努斯}和总督对他的无理和胆大妄为的到来极为不悦，便派遣官员把他赶出城去。官员们来到\Person{路西乌斯}那里，考虑到民众对他非常愤怒且骚动不安，他们不敢自己把他带出房子，以免他被群众打死。他们将此事报告了法官。不久，法官们自己，即公爵\Person{图拉真努斯}和总督\Person{塔提安努斯}，带着许多士兵来到现场，进入房子，在当天的第七时辰，即托特月27日，亲自把\Person{路西乌斯}带了出来。当\Person{路西乌斯}跟着法官们，全城的民众跟在他们后面，无论是基督徒、异教徒还是各种宗教的信徒，都异口同声、同心同德，从他被带出的房子开始，穿过城中心，直到公爵的住宅，不停地呼喊，向他抛掷辱骂和罪名，并高喊：\Quote{把他赶出城去！}然而，公爵把他带进自己家里，\Person{路西乌斯}与他共度了当天剩下的时辰和整个夜晚，在同月28日，公爵一早就押送他远至\Place{尼科波利斯}，把他交给士兵押送出\Place{埃及}。\par

19. \Person{亚大纳削}在帕孔月8日去世，在他安息前5天，他祝圣了\Person{伯多禄}，一位资深的长老，为主教，\Person{伯多禄}继承了他的主教职务，在一切事上效法他。在他之后，他的弟[弟]\Person{弟茂德}继任主教职4年。之后\Person{德奥斐罗}由执事被祝圣为主教(385年)。全书终。\par

\SourceRecord{Translated by Archibald Robertson.；Nicene and Post-Nicene Fathers, Second Series；Vol. 4.；Edited by Philip Schaff and Henry Wace.；Buffalo, NY: Christian Literature Publishing Co.,；1892.；<http://www.newadvent.org/fathers/2820.htm>.}
